% GENERATED FILE. Edit canonical lessons, then run npm run generate:books.
% canonical-source-hash: fnv1a64:ac04074cd1618b21
% canonical-lessons: KA-C312-hogabeku, KA-C312-banda-karana, KA-R312-advice-and-reasons

\chapter{Must, Must Not, and Because}
\label{ch:ka-must-must-not-and-because}
\hlchaptermodality{312}

\begin{chapteropening}
\textbf{By the end of this chapter:} \emph{I can say what I or someone else must or should do with -bēku, what must not be done with -bāradu, and give the reason with … kāraṇa or yākendare.}

Complete the last lesson of chapter 312: give advice with -bēku and -bāradu, and the reasons for it with kāraṇa and yākendare.
\end{chapteropening}

\section[\texorpdfstring{\kn{ಹೋಗಬೇಕು} (hōgabēku)}{ಹೋಗಬೇಕು (hōgabēku)}]{\kn{ಹೋಗಬೇಕು} (hōgabēku) — must, should, must not}

\label{lesson:KA-C312-hogabeku}



\begin{quote}
Say \textquotedblleft{}I want tea\textquotedblright{}. (\emph{nanage cahā bēku}.) Then \textquotedblleft{}Yesterday it rained. So I stayed at home.\textquotedblright{} (\emph{ninne maḷe bantu. hāgāgi nānu maneyallē idde.})

\textbf{\kn{ಬೇಕು}} says a thing is needed. Put it on a verb, and the action is needed.
\end{quote}

\subsection*{You'll want to know: -\kn{ಬೇಕು} on a verb}
\begin{quote}
\textbf{\kn{ನಾನು} \kn{ನಾಳೆ} \kn{ಬೇಗ} \kn{ಹೋಗಬೇಕು}.}
\end{quote}

\begin{quote}
\emph{nānu nāḷe bēga hōgabēku.} — \textquotedblleft{}I must go early tomorrow.\textquotedblright{}
\end{quote}

The verb drops its \textbf{-\kn{ಉ}} and takes \textbf{-\kn{ಅ}}: \textbf{\kn{ಹೋಗು}} $\to$ \textbf{\kn{ಹೋಗ}-}, then \textbf{\kn{ಬೇಕು}}: \textbf{\kn{ಹೋಗಬೇಕು}}. In the same way \textbf{\kn{ಮಾಡಬೇಕು}}, \textbf{\kn{ಓದಬೇಕು}}, \textbf{\kn{ಬರಬೇಕು}}. A verb in \textbf{-\kn{ಇ}} or \textbf{-\kn{ಎ}} adds \textbf{-\kn{ಯ}} first: \textbf{\kn{ಕುಡಿ}} $\to$ \textbf{\kn{ಕುಡಿಯಬೇಕು}}, \textbf{\kn{ತೊಳೆ}} $\to$ \textbf{\kn{ತೊಳೆಯಬೇಕು}}.

\textbf{\kn{ಬೇಕು}} stays the same for every person: \textbf{\kn{ನಾನು} \kn{ಹೋಗಬೇಕು}}, \textbf{\kn{ಅವಳು} \kn{ಹೋಗಬೇಕು}}, \textbf{\kn{ನಾವು} \kn{ಹೋಗಬೇಕು}}. Said about someone else, it is advice:

\begin{quote}
\textbf{\kn{ನೀವು} \kn{ವೈದ್ಯರನ್ನು} \kn{ನೋಡಬೇಕು}.}
\end{quote}

\begin{quote}
\emph{nīvu vaidyarannu nōḍabēku.} — \textquotedblleft{}You should see a doctor.\textquotedblright{}
\end{quote}

\textbf{\kn{ವೈದ್ಯರನ್ನು}} (\emph{vaidyarannu}, \textquotedblleft{}the doctor\textquotedblright{}): \textbf{\kn{ವೈದ್ಯ}} with the respectful \textbf{-\kn{ರು}}, as the one you see.

\subsection*{You'll want to know: -\kn{ಬಾರದು}, must not}
\begin{quote}
\textbf{\kn{ರಾತ್ರಿ} \kn{ತುಂಬಾ} \kn{ಕಾಫಿ} \kn{ಕುಡಿಯಬಾರದು}.}
\end{quote}

\begin{quote}
\emph{rātri tumbā kāphi kuḍiyabāradu.} — \textquotedblleft{}You shouldn't drink a lot of coffee at night.\textquotedblright{}
\end{quote}

\textbf{\kn{ಬಾರದು}} sits where \textbf{\kn{ಬೇಕು}} sat: \textquotedblleft{}must not drink\textquotedblright{} is \textbf{\kn{ಕುಡಿಯಬಾರದು}}; \textquotedblleft{}must not go\textquotedblright{} is \textbf{\kn{ಹೋಗಬಾರದು}}. With no person named, it is a rule for anyone, and the right shape for a notice: \textbf{\kn{ಇಲ್ಲಿ} \kn{ಕಸ} \kn{ಹಾಕಬಾರದು}} (\emph{illi kasa hākabāradu}, \textquotedblleft{}do not throw rubbish here\textquotedblright{}).

\subsection*{Your turn}
Say these aloud:

\begin{itemize}
\raggedright
  \item \textquotedblleft{}I must go early tomorrow\textquotedblright{} — \emph{nānu nāḷe bēga hōgabēku}
  \item \textquotedblleft{}You should see a doctor\textquotedblright{} — \emph{nīvu vaidyarannu nōḍabēku}
  \item \textquotedblleft{}We must buy a ticket\textquotedblright{} — \emph{nāvu ṭikeṭ kharīdisabēku}
  \item \textquotedblleft{}You shouldn't drink a lot of coffee at night\textquotedblright{} — \emph{rātri tumbā kāphi kuḍiyabāradu}
\end{itemize}

\begin{tcolorbox}[breakable,colback=teal!4,colframe=teal!35!black,title={Before you move on}]
\textquotedblleft{}I must go\textquotedblright{}? (\textbf{\kn{ಹೋಗಬೇಕು}}.) \textquotedblleft{}I must drink\textquotedblright{}? (\textbf{\kn{ಕುಡಿಯಬೇಕು}}.) \textquotedblleft{}Must not go\textquotedblright{}? (\textbf{\kn{ಹೋಗಬಾರದು}}.) Does \textbf{\kn{ಬೇಕು}} change for \textquotedblleft{}she\textquotedblright{} or \textquotedblleft{}we\textquotedblright{}? (\textbf{No.}) Which sentence had the same \textbf{\kn{ಬೇಕು}} with no verb? (\textbf{\kn{ನನಗೆ} \kn{ಚಹಾ} \kn{ಬೇಕು}}.)
\end{tcolorbox}
\section[\texorpdfstring{\kn{ಬಂದ} \kn{ಕಾರಣ} (banda kāraṇa)}{ಬಂದ ಕಾರಣ (banda kāraṇa)}]{\kn{ಬಂದ} \kn{ಕಾರಣ} (banda kāraṇa) — giving the reason}

\label{lesson:KA-C312-banda-karana}



\begin{quote}
Say \textquotedblleft{}a reason\textquotedblright{} (\emph{kāraṇa}) and \textquotedblleft{}because\textquotedblright{} (\emph{yākendare}). Then \textquotedblleft{}You should see a doctor.\textquotedblright{} (\emph{nīvu vaidyarannu nōḍabēku}.)
\end{quote}

\subsection*{You'll want to know: the reason first, with \kn{ಕಾರಣ}}
\begin{quote}
\textbf{\kn{ಮಳೆ} \kn{ಬಂದ} \kn{ಕಾರಣ} \kn{ನಾನು} \kn{ಮನೆಯಲ್ಲೇ} \kn{ಇದ್ದೆ}.}
\end{quote}

\begin{quote}
\emph{maḷe banda kāraṇa nānu maneyallē idde.} — \textquotedblleft{}Because it rained, I stayed at home.\textquotedblright{}
\end{quote}

\textbf{\kn{ಬಂದ}} is the past without its ending, the piece inside \textbf{\kn{ಬಂದಾಗ}} and \textbf{\kn{ಮಾಡಿದ} \kn{ನಂತರ}}. With \textbf{\kn{ಕಾರಣ}}, \textquotedblleft{}reason\textquotedblright{}, after it, it reads \textquotedblleft{}(for) the reason (that) it came\textquotedblright{}. Any verb fits:

\begin{quote}
\textbf{\kn{ಅವನಿಗೆ} \kn{ಜ್ವರ} \kn{ಇದ್ದ} \kn{ಕಾರಣ} \kn{ಅವನು} \kn{ಬರಲಿಲ್ಲ}.}
\end{quote}

\begin{quote}
\emph{avanige jvara idda kāraṇa avanu baralilla.} — \textquotedblleft{}Because he had a fever, he did not come.\textquotedblright{}
\end{quote}

\textbf{\kn{ಅವನಿಗೆ}} (\emph{avanige}, \textquotedblleft{}to him\textquotedblright{}) is \textbf{\kn{ಅವನು}} in the \textquotedblleft{}to\textquotedblright{} shape, as \textbf{\kn{ನನಗೆ}} is \textbf{\kn{ನಾನು}}.

\subsection*{You'll want to know: the reason after, with \kn{ಯಾಕೆಂದರೆ}}
\begin{quote}
\textbf{\kn{ನಾನು} \kn{ಇವತ್ತು} \kn{ಮನೆಯಲ್ಲೇ} \kn{ಇರುತ್ತೇನೆ}. \kn{ಯಾಕೆಂದರೆ} \kn{ನನಗೆ} \kn{ಜ್ವರ} \kn{ಇದೆ}.}
\end{quote}

\begin{quote}
\emph{nānu ivattu maneyallē iruttēne. yākendare nanage jvara ide.} — \textquotedblleft{}I'll stay at home today, because I have a fever.\textquotedblright{}
\end{quote}

\textbf{\kn{ಯಾಕೆಂದರೆ}} opens a second sentence, and that sentence is the reason. In writing you will meet its twin, \textbf{\kn{ಏಕೆಂದರೆ}}. So there are two orders, and both are good Kannada: the reason first, ending in \textbf{-\kn{ಅ} \kn{ಕಾರಣ}}, then the result; or the result first, then \textbf{\kn{ಯಾಕೆಂದರೆ}} and the reason.

When someone asks you why, answer in either order:

\begin{quote}
\textbf{\kn{ನೀನು} \kn{ನಿನ್ನೆ} \kn{ಯಾಕೆ} \kn{ಬರಲಿಲ್ಲ}?}
\end{quote}

\begin{quote}
\emph{nīnu ninne yāke baralilla?} — \textquotedblleft{}Why didn't you come yesterday?\textquotedblright{}
\end{quote}

\subsection*{Your turn}
Say these aloud:

\begin{itemize}
\raggedright
  \item \textquotedblleft{}Because it rained, I stayed at home\textquotedblright{} — \emph{maḷe banda kāraṇa nānu maneyallē idde}
  \item \textquotedblleft{}Because he had a fever, he did not come\textquotedblright{} — \emph{avanige jvara idda kāraṇa avanu baralilla}
  \item \textquotedblleft{}I'll stay at home today, because I have a fever\textquotedblright{} — \emph{nānu ivattu maneyallē iruttēne. yākendare nanage jvara ide.}
  \item \textquotedblleft{}Why didn't you come yesterday?\textquotedblright{} — \emph{nīnu ninne yāke baralilla?}
\end{itemize}

\begin{tcolorbox}[breakable,colback=teal!4,colframe=teal!35!black,title={Before you move on}]
\textquotedblleft{}Because it rained\textquotedblright{}, reason first? (\textbf{\kn{ಮಳೆ} \kn{ಬಂದ} \kn{ಕಾರಣ}}.) Which word opens a reason that comes after? (\textbf{\kn{ಯಾಕೆಂದರೆ}}, or \textbf{\kn{ಏಕೆಂದರೆ}} in writing.) \textquotedblleft{}Why didn't you come?\textquotedblright{} (\textbf{\kn{ಯಾಕೆ} \kn{ಬರಲಿಲ್ಲ}?})
\end{tcolorbox}
\section[(review)]{Review — advice and reasons}

\label{lesson:KA-R312-advice-and-reasons}



\begin{quote}
Read \textbf{\kn{ಸ್ವಚ್ಛ}} aloud, and name its small letter. (\emph{svaccha}; \textbf{\kn{ಛ}}.) Then \textquotedblleft{}So I stayed at home.\textquotedblright{} (\emph{hāgāgi nānu maneyallē idde}.)
\end{quote}

\subsection*{Your turn}
Say these aloud:

\begin{itemize}
\raggedright
  \item You have a fever. You should see a doctor.
  \item You must not go to school today.
  \item Because the bus did not come, I walked.
  \item I must go early tomorrow, because Grandmother is coming.
  \item I think we must clean the house.
  \item The bus didn't come. That's why I came late.
\end{itemize}

\begin{tcolorbox}[breakable,colback=teal!4,colframe=teal!35!black,title={Before you move on}]
\textquotedblleft{}Must go\textquotedblright{}? (\textbf{\kn{ಹೋಗಬೇಕು}}.) \textquotedblleft{}Must not go\textquotedblright{}? (\textbf{\kn{ಹೋಗಬಾರದು}}.) \textquotedblleft{}Because it came\textquotedblright{}, reason first? (\textbf{\kn{ಬಂದ} \kn{ಕಾರಣ}}.) The reason after? (\textbf{\kn{ಯಾಕೆಂದರೆ}} …)
\end{tcolorbox}
